% !TeX program = xelatex
\documentclass[11pt,a4paper,fontset=fandol]{ctexart}

\usepackage{geometry}
\geometry{left=2.5cm, right=2.5cm, top=2.5cm, bottom=2.5cm}

\usepackage{amsmath, amssymb, amsthm}
\usepackage{mathtools}
\usepackage{bm}
\allowdisplaybreaks

\newtheorem{theorem}{定理}[section]
\newtheorem{lemma}[theorem]{引理}
\newtheorem{property}[theorem]{性质}
\newtheorem{corollary}[theorem]{推论}
\newtheorem{example}{例题}[section]

\theoremstyle{remark}
\newtheorem{remark}{注}[section]

\usepackage{newtxtext}
\usepackage{newtxmath}
\usepackage[most]{tcolorbox}
\newtcbtheorem[number within=section]{definition}{定义}{
    enhanced,
    boxrule=0.5pt,
    arc=0pt,
    colback=gray!4!white,
    colframe=black!80,
    fonttitle=\bfseries\sffamily,
    coltitle=black,
    attach boxed title to top left={yshift=-2mm, xshift=2mm},
    boxed title style={boxrule=0pt, colframe=white, colback=white},
    separator sign={\ },
    description font=\normalfont
}{def}

\usepackage{booktabs}
\usepackage{graphicx}
\usepackage{enumitem}
\setlist{nosep, leftmargin=2em}
\usepackage{caption}
\captionsetup{font=small, labelfont=bf, justification=centering}
\usepackage{tikz}
\usetikzlibrary{arrows.meta, decorations.markings, calc, patterns}

\usepackage{hyperref}
\usepackage{cleveref}
\hypersetup{
    colorlinks=true,
    linkcolor=black,
    citecolor=black,
    urlcolor=black
}

\ctexset{
    section/name={,},
    section/number=\arabic{section},
    section/format=\Large\bfseries,
    subsection/number=\arabic{section}.\arabic{subsection},
    subsection/format=\large\bfseries,
    subsubsection/number=\arabic{section}.\arabic{subsection}.\arabic{subsubsection},
    subsubsection/format=\normalsize\bfseries
}

\renewcommand{\baselinestretch}{1.15}
\setlength{\parskip}{0.3em}
\setlength{\parindent}{2em}

\usepackage{fancyhdr}
\pagestyle{fancy}
\fancyhf{}
\fancyhead[L]{\small 概率论与数理统计}
\fancyhead[R]{\small \leftmark}
\fancyfoot[C]{\small\thepage}
\renewcommand{\headrulewidth}{0.4pt}

\begin{document}

\pagenumbering{gobble}
\title{概率论与数理统计笔记}
\author{}
\date{}
\maketitle
\tableofcontents
\clearpage
\pagenumbering{arabic}

\section{概率论基本概念}
\subsection{随机试验、样本空间与随机事件}
概率论以随机现象的统计规律性为研究对象。随机试验可在相同条件下重复进行，所有可能结果可确定，而单次试验结果不能预先确定。

\begin{definition}{随机试验与样本空间}{random_experiment}
在相同条件下可以重复进行、具有多个可确定的可能结果且单次结果不能预先确定的试验，称为\textbf{随机试验}。随机试验的所有可能结果构成的集合称为\textbf{样本空间}，记为 $\Omega$；其中的每个结果称为\textbf{样本点}。
\end{definition}

样本空间由试验及其观察目标确定。例如，连续抛掷两次硬币并记录先后次序时，
\[
    \Omega=\{(H,H),(H,T),(T,H),(T,T)\}.
\]
若观察量为正面出现的次数，则样本空间为 $\{0,1,2\}$。

\begin{definition}{随机事件}{random_event}
样本空间 $\Omega$ 的子集称为\textbf{随机事件}。全体事件记为 $\mathcal{F}$；在概率模型中，$\mathcal{F}$ 对补集和可列并封闭，并包含 $\Omega$。若一次试验的结果属于事件 $A$，就称事件 $A$ 发生。$\Omega$ 称为必然事件，$\varnothing$ 称为不可能事件，只含一个样本点的事件称为基本事件。
\end{definition}

事件是集合，因而可以用集合关系和集合运算描述。例如，$A\subseteq B$ 表示 $A$ 发生必然导致 $B$ 发生；$A\cup B$ 表示至少一个发生；$A\cap B$ 表示两者同时发生；$A^c=\Omega\setminus A$ 表示 $A$ 不发生。

\subsection{事件关系与事件运算}
事件关系与运算如下：
\begin{itemize}
    \item $A\subseteq B$：$A$ 蕴含 $B$；$A=B$：两个事件相等。
    \item $A\cap B=\varnothing$：$A$ 与 $B$ 互不相容（互斥）。
    \item $A\cup B$：$A$ 与 $B$ 至少一个发生；$A\cap B$：$A$ 与 $B$ 同时发生。
    \item $A\setminus B=A\cap B^c$：$A$ 发生而 $B$ 不发生。
\end{itemize}
事件运算遵循集合运算律，其中德摩根律为
\[
    (A\cup B)^c=A^c\cap B^c,\qquad
    (A\cap B)^c=A^c\cup B^c.
\]
复杂事件可表示为集合运算的组合，并由事件运算律及概率公式求其概率。

\subsection{频率、概率公理与概率性质}
在 $n$ 次重复试验中，若事件 $A$ 发生了 $N_n(A)$ 次，则
\[
    f_n(A)=\frac{N_n(A)}{n}
\]
称为事件 $A$ 的频率。若各次试验相互独立且事件 $A$ 在每次试验中的概率均为 $p$，则由大数定律，$f_n(A)$ 依概率收敛于 $p$。有限次试验所得频率是相对次数，不必等于 $p$。

\begin{definition}{概率}{probability}
概率是定义在事件集合 $\mathcal{F}$ 上的数值函数 $P:\mathcal{F}\to[0,1]$，满足：
\begin{enumerate}
    \item 非负性：对任意事件 $A$，$P(A)\geq 0$；
    \item 规范性：$P(\Omega)=1$；
    \item 可列可加性：若 $A_1,A_2,\ldots\in\mathcal{F}$ 两两互斥，则
    \[
        P\left(\bigcup_{i=1}^{\infty}A_i\right)
        =\sum_{i=1}^{\infty}P(A_i).
    \]
\end{enumerate}
\end{definition}

由概率公理可推出以下性质：
\begin{property}[概率的基本运算性质]
\[
\begin{aligned}
    &P(\varnothing)=0,\qquad P(A^c)=1-P(A),\\
    &A\subseteq B\ \Longrightarrow\ P(A)\leq P(B),\\
    &P(A\setminus B)=P(A)-P(A\cap B),\\
    &P(A\cup B)=P(A)+P(B)-P(A\cap B).
\end{aligned}
\]
\end{property}
若 $A$ 与 $B$ 互斥，则 $P(A\cup B)=P(A)+P(B)$。只有在互斥时才能直接把并事件概率写成概率之和。

\subsection{古典概型与几何概型}
当随机试验只有有限个可能结果，且每个基本结果出现的可能性相同时，称为\textbf{古典概型}。若样本空间含 $n$ 个等可能基本结果，事件 $A$ 含 $m$ 个基本结果，则
\[
    P(A)=\frac{m}{n}.
\]
该公式的条件是样本点等可能；若样本点不等可能，则不能用有利样本点数与样本点总数之比计算概率。

\begin{example}[掷均匀骰子]
掷一颗均匀骰子，求出现偶数点的概率。

\textbf{解：}样本空间为 $\{1,2,3,4,5,6\}$，六个结果等可能；偶数点事件为 $\{2,4,6\}$，所以
\[
    P(\text{偶数点})=\frac{3}{6}=\frac{1}{2}.
\]
\end{example}

若试验结果均匀地落在区域 $D$ 内，且所关心的事件对应区域 $G\subseteq D$，则称为\textbf{几何概型}，其概率为
\[
    P(G)=\frac{\mu(G)}{\mu(D)},
\]
其中 $\mu$ 表示相应维数的长度、面积或体积测度。该公式成立的条件是试验结果在 $D$ 上关于 $\mu$ 均匀分布。

\subsection{条件概率与乘法公式}
当已知事件 $B$ 已发生时，样本空间缩小为 $B$，再计算 $A$ 发生的概率，这就是条件概率。

\begin{definition}{条件概率}{conditional_probability}
若 $P(B)>0$，则在事件 $B$ 已发生的条件下事件 $A$ 的概率定义为
\[
    P(A\mid B)=\frac{P(A\cap B)}{P(B)}.
\]
\end{definition}

固定 $B$ 且 $P(B)>0$ 时，$P(\,\cdot\mid B)$ 本身也是一个概率。由定义立即得到乘法公式：
\[
    P(A\cap B)=P(B)P(A\mid B)=P(A)P(B\mid A),
\]
相应的条件概率有定义时，两种写法均可使用。多个事件的乘法公式按顺序展开：
\[
    P\left(\bigcap_{i=1}^{n}A_i\right)
    =P(A_1)\prod_{k=2}^{n}
    P\left(A_k\mid\bigcap_{i=1}^{k-1}A_i\right).
\]
该式在右侧出现的条件概率均有定义时成立；若某个前缀交事件的概率为 $0$，则完整交事件的概率也为 $0$。
乘法公式用于“同时发生”的概率；若各次试验独立，条件概率才可以进一步简化。

\subsection{全概率公式与贝叶斯公式}
设 $A_1,\ldots,A_n$ 两两互斥且 $\bigcup_{i=1}^n A_i=\Omega$，并且 $P(A_i)>0$。这组事件构成样本空间的一个划分。对任意事件 $B$，
\[
    P(B)=\sum_{i=1}^{n}P(A_i)P(B\mid A_i).
\]
该式将 $B$ 表示为各互斥事件 $A_i$ 下发生概率的加权和，权重为 $P(A_i)$。

\begin{theorem}[贝叶斯公式]
在上述划分下，若 $P(B)>0$，则
\[
    P(A_k\mid B)
    =\frac{P(A_k)P(B\mid A_k)}
    {\displaystyle\sum_{i=1}^{n}P(A_i)P(B\mid A_i)}.
\]
该式给出条件概率 $P(A_k\mid B)$。
\end{theorem}

\begin{example}[从不同箱子中抽球]
等可能地选取一个箱子：1 号箱有 1 个红球和 4 个白球，2 号箱有 2 个红球和 3 个白球，3 号箱有 3 个红球。已知抽出的是红球，求球来自 1 号箱的概率。

\textbf{解：}令 $A_i$ 表示选中 $i$ 号箱，$R$ 表示抽到红球。则
\[
    P(A_1)=P(A_2)=P(A_3)=\frac{1}{3},\quad
    P(R\mid A_1)=\frac{1}{5},\quad
    P(R\mid A_2)=\frac{2}{5},\quad
    P(R\mid A_3)=1.
\]
先用全概率公式求
\[
    P(R)=\frac{1}{3}\left(\frac{1}{5}+\frac{2}{5}+1\right)=\frac{8}{15}.
\]
再用贝叶斯公式：
\[
    P(A_1\mid R)
    =\frac{P(A_1)P(R\mid A_1)}{P(R)}
    =\frac{(1/3)(1/5)}{8/15}
    =\frac{1}{8}.
\]
\end{example}

\subsection{事件的独立性}
\begin{definition}{两个事件的独立性}{event_independence}
若事件 $A$ 与 $B$ 满足
\[
    P(A\cap B)=P(A)P(B),
\]
则称 $A$ 与 $B$ 相互独立。若 $P(B)>0$，这等价于 $P(A\mid B)=P(A)$。
\end{definition}

事件 $A_1,\ldots,A_n$ 相互独立，当且仅当对任意非空指标子集 $I\subseteq\{1,\ldots,n\}$，均有
\[
    P\left(\bigcap_{i\in I}A_i\right)=\prod_{i\in I}P(A_i).
\]
当 $n=3$ 时，相互独立要求对每一对不同指标 $i,j\in\{1,2,3\}$，
\[
    P(A_i\cap A_j)=P(A_i)P(A_j),
\]
并且
\[
    P(A_1\cap A_2\cap A_3)=P(A_1)P(A_2)P(A_3).
\]
两两独立只约束各事件对，不蕴含三事件相互独立。

\begin{example}[两次独立掷硬币]
均匀硬币连续掷两次，令 $A$ 为“第一次正面”，$B$ 为“第二次正面”。由于四种有序结果等可能，
\[
    P(A)=P(B)=\frac{1}{2},\qquad
    P(A\cap B)=\frac{1}{4}=P(A)P(B),
\]
所以 $A$ 与 $B$ 独立。
\end{example}

\begin{remark}[互斥与独立]
互斥与独立是不同的事件关系。若 $A$ 与 $B$ 互斥且 $P(A)>0$、$P(B)>0$，则 $P(A\cap B)=0\ne P(A)P(B)$，故二者不独立。若 $P(B)>0$，独立性等价于 $P(A\mid B)=P(A)$。
\end{remark}

\clearpage
\section{随机变量及其分布}
\subsection{随机变量与分布函数}
随机变量把随机试验结果映射为实数，使事件可以用不等式或方程表示。例如，掷骰子时可令 $X$ 表示点数；三次投币时可令 $X$ 表示正面出现次数。

\begin{definition}{随机变量}{random_variable}
设概率空间为 $(\Omega,\mathcal{F},P)$。若映射 $X:\Omega\to\mathbb{R}$ 满足对任意 Borel 集 $B\subseteq\mathbb{R}$，原像 $\{\omega\in\Omega:X(\omega)\in B\}$ 属于 $\mathcal{F}$，则称 $X$ 为随机变量。
\end{definition}

离散型随机变量的可能取值有限或可列；连续型随机变量的分布函数可由概率密度积分得到。随机变量还可具有既含离散概率质量又含连续密度部分的混合分布；本章分别讨论离散型和连续型情形。

\begin{definition}{分布函数}{distribution_function}
随机变量 $X$ 的分布函数定义为
\[
    F_X(x)=P(X\leq x),\qquad -\infty<x<\infty.
\]
\end{definition}

分布函数完整刻画随机变量的概率分布，并满足：
\begin{enumerate}
    \item 单调不减：若 $x_1<x_2$，则 $F_X(x_1)\leq F_X(x_2)$；
    \item 有界：$0\leq F_X(x)\leq 1$；
    \item 极限：$\lim_{x\to-\infty}F_X(x)=0$，$\lim_{x\to+\infty}F_X(x)=1$；
    \item 右连续：$\lim_{t\downarrow x}F_X(t)=F_X(x)$。
\end{enumerate}
任意 $a<b$，有
\[
    P(a<X\leq b)=F_X(b)-F_X(a).
\]
若 $F_X$ 在 $x$ 处连续，则 $P(X=x)=0$；若有跳跃，则
\[
    P(X=x)=F_X(x)-F_X(x-),
\]
其中 $F_X(x-)$ 是左极限。

\subsection{离散型随机变量及其分布律}
若随机变量 $X$ 的所有可能取值为 $x_1,x_2,\ldots$，并记
\[
    p_k=P(X=x_k),
\]
则称 $\{p_k\}$ 为 $X$ 的分布律（概率质量函数）。它满足
\[
    p_k\geq 0,\qquad \sum_k p_k=1.
\]
离散型随机变量的分布函数可由分布律累加得到：
\[
    F_X(x)=\sum_{x_k\leq x}p_k.
\]

\begin{example}[由分布律计算区间概率]
设 $X$ 的分布律为 $P(X=0)=0.2$、$P(X=1)=0.5$、$P(X=2)=0.3$。则
\[
    P(X\leq 1)=0.2+0.5=0.7,\qquad
    P(0<X\leq 2)=0.5+0.3=0.8.
\]
\end{example}

\subsection{两点分布、二项分布与泊松分布}
常用的离散型分布包括：
\begin{itemize}
    \item \textbf{两点分布（$0$--$1$ 分布）}：若 $X$ 只取 $0,1$，且 $P(X=1)=p$，$P(X=0)=1-p$，记为 $X\sim\operatorname{Bernoulli}(p)$。它表示仅有两种可能结果的试验。
    \item \textbf{二项分布}：在 $n$ 次相互独立、成功概率均为 $p$ 的伯努利试验中，令 $X$ 为成功次数，则
    \[
        P(X=k)=\binom{n}{k}p^k(1-p)^{n-k},
        \quad k=0,1,\ldots,n,
        \qquad X\sim B(n,p).
    \]
    \item \textbf{泊松分布}：若
    \[
        P(X=k)=e^{-\lambda}\frac{\lambda^k}{k!},
        \quad k=0,1,2,\ldots,\quad \lambda>0,
    \]
    则称 $X$ 服从参数为 $\lambda$ 的泊松分布，记为 $X\sim P(\lambda)$。在泊松过程中，长度为 $t$ 的区间内事件数服从参数为 $\lambda t$ 的泊松分布，其中 $\lambda$ 为单位区间的发生率。
\end{itemize}

\begin{example}[二项分布的概率]
独立投掷一枚均匀硬币 4 次，令 $X$ 为正面次数。则 $X\sim B(4,1/2)$，恰有 2 次正面的概率为
\[
    P(X=2)=\binom{4}{2}\left(\frac12\right)^2\left(\frac12\right)^2
    =\frac{3}{8}.
\]
\end{example}

当 $n$ 较大、$p$ 较小且 $np=\lambda$ 适中时，二项分布可用参数 $\lambda=np$ 的泊松分布近似：
\[
    \binom{n}{k}p^k(1-p)^{n-k}
    \approx e^{-np}\frac{(np)^k}{k!}.
\]
这是近似而非恒等式；使用时应回到具体问题判断 $n,p$ 是否适合。

\subsection{连续型随机变量与概率密度}
若存在非负可积函数 $f_X(x)$，使得对任意实数 $x$，
\[
    F_X(x)=\int_{-\infty}^{x}f_X(t)\,\mathrm{d}t,
\]
则称 $X$ 为连续型随机变量，$f_X$ 为其概率密度函数。密度满足
\[
    f_X(x)\geq 0,\qquad
    \int_{-\infty}^{+\infty}f_X(x)\,\mathrm{d}x=1.
\]
对 $a<b$，
\[
    P(a<X\leq b)=\int_a^b f_X(x)\,\mathrm{d}x.
\]
若 $F_X$ 在 $x$ 处可导，则 $f_X(x)=F_X'(x)$（在相应的几乎处处意义下）。密度在单点的函数值不是该点的概率；对连续型随机变量，$P(X=x)=0$。

\begin{example}[确定密度常数并计算概率]
设
\[
    f_X(x)=
    \begin{cases}
        c x, & 0<x<2,\\
        0, & \text{其他}.
    \end{cases}
\]
由密度积分为 1 得 $1=\int_0^2 cx\,\mathrm{d}x=2c$，故 $c=1/2$。因此
\[
    P(1<X<2)=\int_1^2\frac{x}{2}\,\mathrm{d}x=\frac{3}{4}.
\]
\end{example}

\subsection{均匀分布、指数分布与正态分布}
\begin{itemize}
    \item \textbf{均匀分布}：若 $a<b$，且
    \[
        f_X(x)=
        \begin{cases}
            \dfrac{1}{b-a}, & a<x<b,\\
            0, & \text{其他},
        \end{cases}
    \]
    则记 $X\sim U(a,b)$。区间内等长子区间具有相同概率。
    \item \textbf{指数分布}：若 $\lambda>0$，且
    \[
        f_X(x)=
        \begin{cases}
            \lambda e^{-\lambda x}, & x>0,\\
            0, & x\leq 0,
        \end{cases}
    \]
    则记 $X\sim\operatorname{Exp}(\lambda)$。其分布函数为 $F_X(x)=1-e^{-\lambda x}$（$x>0$），而 $x\leq0$ 时为 $0$。
    \item \textbf{正态分布}：若 $\mu\in\mathbb{R}$、$\sigma>0$，且
    \[
        f_X(x)=\frac{1}{\sqrt{2\pi}\sigma}
        \exp\left[-\frac{(x-\mu)^2}{2\sigma^2}\right],
        \qquad -\infty<x<\infty,
    \]
    则记 $X\sim N(\mu,\sigma^2)$。密度曲线关于 $\mu$ 对称，呈钟形；$\mu$ 为位置参数，$\sigma$ 决定离散程度。
\end{itemize}

若 $X\sim N(\mu,\sigma^2)$，标准化变量
\[
    Z=\frac{X-\mu}{\sigma}
\]
服从标准正态分布 $N(0,1)$。记标准正态分布函数为
\[
    \Phi(z)=\frac{1}{\sqrt{2\pi}}\int_{-\infty}^{z}e^{-t^2/2}\,\mathrm{d}t.
\]
于是对 $a<b$，
\[
    P(a<X\leq b)
    =\Phi\left(\frac{b-\mu}{\sigma}\right)
    -\Phi\left(\frac{a-\mu}{\sigma}\right),
    \qquad \Phi(-z)=1-\Phi(z).
\]

\begin{example}[正态分布标准化]
若 $X\sim N(10,4)$，则 $\sigma=2$，从而
\[
    P(8<X\leq 12)
    =\Phi(1)-\Phi(-1)
    =2\Phi(1)-1.
\]
\end{example}

\subsection{随机变量函数的分布}
若 $Y=g(X)$，则 $Y$ 也是随机变量。其分布函数可由下式求得：
\[
    F_Y(y)=P(Y\leq y)=P\bigl(g(X)\leq y\bigr).
\]
先把事件 $\{g(X)\leq y\}$ 化为关于 $X$ 的取值范围，再利用 $X$ 的分布求概率。

\begin{example}[离散型随机变量的函数]
若 $X$ 取值 $-1,0,1$ 的概率分别为 $0.2,0.5,0.3$，令 $Y=X^2$。则 $Y$ 取值 $0,1$，并且
\[
    P(Y=0)=P(X=0)=0.5,\qquad
    P(Y=1)=P(X=-1)+P(X=1)=0.5.
\]
求离散型函数的分布时，要合并所有映射到同一函数值的原取值及其概率。
\end{example}

若 $X$ 为连续型随机变量，且 $Y=g(X)$ 在取值范围内严格单调、可微并存在反函数 $x=g^{-1}(y)$，则变量变换公式为
\[
    f_Y(y)=f_X\bigl(g^{-1}(y)\bigr)
    \left|\frac{\mathrm{d}}{\mathrm{d}y}g^{-1}(y)\right|,
\]
其中 $y$ 只取 $g(X)$ 的实际取值范围，范围外密度为 $0$。若 $g$ 非单调，应把所有满足 $g(x_i)=y$ 的原像分支贡献相加：
\[
    f_Y(y)=\sum_{x_i:\,g(x_i)=y}
    \frac{f_X(x_i)}{|g'(x_i)|},
\]
适用于各分支可逆且导数非零的情形。

\begin{example}[正数的线性变换]
设 $X$ 的密度为 $f_X$，令 $Y=2X+3$。变换单调递增，$x=(y-3)/2$ 且 $\mathrm{d}x/\mathrm{d}y=1/2$，故
\[
    f_Y(y)=\frac{1}{2}f_X\left(\frac{y-3}{2}\right),
\]
其定义范围是 $y=2x+3$ 对 $x$ 的取值范围所对应的区间。
\end{example}

\begin{remark}[分布与变量变换的条件]
分布律中各概率非负且总和为 1；连续型概率由密度积分给出；分布函数以 $P(X\leq x)$ 定义并右连续；正态分布 $N(\mu,\sigma^2)$ 的第二个参数为方差。变量变换密度含反函数导数的绝对值；非单调变换的各可逆分支均对密度作出贡献。
\end{remark}

\clearpage
\section{多维随机变量及其分布}
\subsection{二维随机变量与联合分布}
两个随机变量共同构成随机向量 $(X,Y)$。其联合分布描述 $X$ 与 $Y$ 取值的概率规律。

\begin{definition}{二维随机变量的联合分布函数}{joint_distribution_function}
二维随机变量 $(X,Y)$ 的联合分布函数定义为
\[
    F_{X,Y}(x,y)=P(X\leq x,\;Y\leq y),
    \qquad x,y\in\mathbb{R}.
\]
\end{definition}

若 $X,Y$ 均为离散型随机变量，记
\[
    p_{ij}=P(X=x_i,Y=y_j),
\]
则 $\{p_{ij}\}$ 为联合分布律，满足
\[
    p_{ij}\geq 0,\qquad \sum_i\sum_jp_{ij}=1.
\]
对任意指标集合 $I,J$，
\[
    P(X\in\{x_i:i\in I\},\,Y\in\{y_j:j\in J\})
    =\sum_{i\in I}\sum_{j\in J}p_{ij}.
\]

若存在非负函数 $f_{X,Y}$，使得对平面上任意可测区域 $D$，
\[
    P\bigl((X,Y)\in D\bigr)
    =\iint_D f_{X,Y}(x,y)\,\mathrm{d}x\,\mathrm{d}y,
\]
则称 $(X,Y)$ 为连续型随机向量，$f_{X,Y}$ 为联合概率密度，且
\[
    f_{X,Y}(x,y)\geq 0,\qquad
    \int_{-\infty}^{+\infty}\int_{-\infty}^{+\infty}
    f_{X,Y}(x,y)\,\mathrm{d}x\,\mathrm{d}y=1.
\]
联合分布函数与密度之间的关系为
\[
    F_{X,Y}(x,y)=\int_{-\infty}^{x}\int_{-\infty}^{y}
    f_{X,Y}(u,v)\,\mathrm{d}v\,\mathrm{d}u.
\]
在密度连续的点，$f_{X,Y}(x,y)=\dfrac{\partial^2}{\partial x\,\partial y}F_{X,Y}(x,y)$。

\subsection{边缘分布}
边缘分布分别描述随机向量各分量的分布。离散型情形下，对联合分布律求和：
\[
    p_X(x_i)=P(X=x_i)=\sum_jp_{ij},\qquad
    p_Y(y_j)=P(Y=y_j)=\sum_ip_{ij}.
\]
连续型情形下，对联合密度沿另一变量积分：
\[
    f_X(x)=\int_{-\infty}^{+\infty}f_{X,Y}(x,y)\,\mathrm{d}y,\qquad
    f_Y(y)=\int_{-\infty}^{+\infty}f_{X,Y}(x,y)\,\mathrm{d}x.
\]
边缘分布函数亦可由联合分布函数得到：
\[
    F_X(x)=\lim_{y\to+\infty}F_{X,Y}(x,y),\qquad
    F_Y(y)=\lim_{x\to+\infty}F_{X,Y}(x,y).
\]
由联合分布可确定边缘分布；一般而言，仅知道边缘分布不能唯一确定联合分布。

\begin{example}[由离散联合分布求边缘分布]
设 $(X,Y)$ 的联合分布律如下：
\[
\begin{array}{c|cc|c}
    &Y=0&Y=1&P_X(x)\\
    \hline
    X=0&\frac12&\frac16&\frac23\\
    X=1&\frac16&\frac16&\frac13\\
    \hline
    P_Y(y)&\frac23&\frac13&1
\end{array}
\]
例如 $P(X=0)=P(X=0,Y=0)+P(X=0,Y=1)=\frac12+\frac16=\frac23$；表中列和给出 $Y$ 的边缘分布。
\end{example}

\begin{example}[由连续联合密度求边缘密度]
设
\[
    f_{X,Y}(x,y)=
    \begin{cases}
        2,&0<x<y<1,\\
        0,&\text{其他}.
    \end{cases}
\]
其关于 $X$ 的边缘密度为
\[
    f_X(x)=\int_x^1 2\,\mathrm{d}y=2(1-x),
    \qquad 0<x<1,
\]
关于 $Y$ 的边缘密度为
\[
    f_Y(y)=\int_0^y 2\,\mathrm{d}x=2y,
    \qquad 0<y<1.
\]
\end{example}

\subsection{条件分布}
离散型情形下，若 $P(Y=y_j)>0$，则
\[
    P(X=x_i\mid Y=y_j)
    =\frac{P(X=x_i,Y=y_j)}{P(Y=y_j)}
    =\frac{p_{ij}}{p_Y(y_j)}.
\]
对固定的 $y_j$，这些条件概率关于 $i$ 的和为 $1$，构成 $Y=y_j$ 条件下 $X$ 的条件分布律。

连续型情形下，若 $f_Y(y)>0$，则 $X$ 在 $Y=y$ 条件下的条件密度定义为
\[
    f_{X\mid Y}(x\mid y)=\frac{f_{X,Y}(x,y)}{f_Y(y)}.
\]
固定 $y$ 后，$f_{X\mid Y}(x\mid y)$ 关于 $x$ 的积分为 $1$。对应的条件分布函数为
\[
    F_{X\mid Y}(x\mid y)=\int_{-\infty}^{x}f_{X\mid Y}(u\mid y)\,\mathrm{d}u.
\]
条件密度公式对满足 $f_Y(y)>0$ 的 $y$ 成立；在 $f_Y(y)=0$ 处不能直接使用该比值定义。

\begin{example}[由联合密度求条件密度]
对前例的密度，在 $0<y<1$ 时，
\[
    f_{X\mid Y}(x\mid y)=
    \begin{cases}
        \dfrac{2}{2y}=\dfrac1y,&0<x<y,\\
        0,&\text{其他}.
    \end{cases}
\]
对任意 $0<y<1$，有 $\int_0^y(1/y)\,\mathrm{d}x=1$。
\end{example}

\subsection{随机变量的相互独立性}
随机变量 $X$ 与 $Y$ 相互独立，当且仅当对任意 $x,y\in\mathbb{R}$，
\[
    F_{X,Y}(x,y)=F_X(x)F_Y(y).
\]
等价地：
\begin{itemize}
    \item 离散型情形：$p_{ij}=p_X(x_i)p_Y(y_j)$，对所有可能取值成立；
    \item 连续型情形：$f_{X,Y}(x,y)=f_X(x)f_Y(y)$，几乎处处成立。
\end{itemize}
此外，若某个条件概率有定义，独立性意味着条件分布等于相应的边缘分布。联合分布可分解为边缘分布乘积，是独立性的判定条件；仅比较期望或协方差一般不足以判定独立。

\begin{example}[用联合分布律判定独立性]
在前述离散联合分布表中，
\[
    P(X=0,Y=0)=\frac12,\qquad
    P(X=0)P(Y=0)=\frac23\cdot\frac23=\frac49.
\]
两者不相等，因此 $X$ 与 $Y$ 不独立。相应地，
\[
    P(X=1\mid Y=0)=\frac{1/6}{2/3}=\frac14
    \ne P(X=1)=\frac13.
\]
\end{example}

\subsection{两个随机变量函数的分布}
令 $Z=X+Y$。当 $X,Y$ 为独立离散型随机变量时，
\[
    P(Z=z)=\sum_{x_i+y_j=z}p_X(x_i)p_Y(y_j).
\]
若 $X,Y$ 独立且具有连续密度，则其和的密度为卷积：
\[
    f_Z(z)=\int_{-\infty}^{+\infty}f_X(x)f_Y(z-x)\,\mathrm{d}x.
\]
独立性是上述乘积形式所需的条件；若 $X,Y$ 不独立，应由联合分布计算和的分布。

设 $M=\max(X,Y)$、$N=\min(X,Y)$，且 $X,Y$ 相互独立，分布函数分别为 $F_X,F_Y$，则
\[
\begin{aligned}
    F_M(z)
    &=P(X\leq z,Y\leq z)
      =F_X(z)F_Y(z),\\
    F_N(z)
    &=1-P(X>z,Y>z)
      =F_X(z)+F_Y(z)-F_X(z)F_Y(z).
\end{aligned}
\]

\begin{example}[独立均匀随机变量的最大值]
设 $X,Y$ 相互独立且均服从 $U(0,1)$。当 $0\leq z\leq1$ 时，
\[
    F_{\max(X,Y)}(z)=F_X(z)F_Y(z)=z^2.
\]
因此最大值的密度为 $f_{\max(X,Y)}(z)=2z$（$0<z<1$），区间外密度为 $0$。
\end{example}

\begin{remark}[独立性条件]
边缘分布不能单独确定联合分布。由边缘分布相乘得到联合分布，仅在随机变量独立时成立；使用卷积公式及最大值、最小值的乘积分布式时，亦须满足独立条件。
\end{remark}

\clearpage
\section{随机变量的数字特征}
\subsection{数学期望}
数学期望刻画随机变量取值的概率加权平均。离散型随机变量 $X$ 的可能取值为 $x_k$，对应概率为 $p_k$。若级数绝对收敛，即
\[
    \sum_k |x_k|p_k<\infty,
\]
则 $X$ 的数学期望为
\[
    E(X)=\sum_k x_kp_k.
\]
连续型随机变量 $X$ 的概率密度为 $f_X$，若
\[
    \int_{-\infty}^{+\infty}|x|f_X(x)\,\mathrm{d}x<\infty,
\]
则
\[
    E(X)=\int_{-\infty}^{+\infty}x f_X(x)\,\mathrm{d}x.
\]
绝对可积条件保证数学期望有限且定义良好。

\begin{example}[离散型随机变量的期望]
设 $X$ 的分布律为 $P(X=0)=0.2$、$P(X=1)=0.5$、$P(X=2)=0.3$，则
\[
    E(X)=0(0.2)+1(0.5)+2(0.3)=1.1.
\]
\end{example}

\subsection{随机变量函数的数学期望}
若 $Y=g(X)$，在相应期望存在时，不必先求 $Y$ 的分布即可计算其期望：
\[
    E[g(X)]=
    \begin{cases}
        \displaystyle\sum_k g(x_k)p_k,&X\text{ 为离散型},\\[2mm]
        \displaystyle\int_{-\infty}^{+\infty}g(x)f_X(x)\,\mathrm{d}x,
        &X\text{ 为连续型}.
    \end{cases}
\]
对二维离散型随机向量，
\[
    E[g(X,Y)]=\sum_i\sum_j g(x_i,y_j)p_{ij};
\]
对具有联合密度的连续型随机向量，
\[
    E[g(X,Y)]=\iint_{\mathbb{R}^2}g(x,y)f_{X,Y}(x,y)\,\mathrm{d}x\,\mathrm{d}y,
\]
以上公式均要求相应绝对期望有限。

\begin{property}[数学期望的性质]
在各期望存在时，
\[
\begin{aligned}
    &E(c)=c,\qquad E(aX+b)=aE(X)+b,\\
    &E(X+Y)=E(X)+E(Y),\\
    &X,Y\text{ 独立}\ \Longrightarrow\ E(XY)=E(X)E(Y).
\end{aligned}
\]
期望的线性性不要求随机变量相互独立；乘积期望分解的充分条件之一是独立性，反向推论一般不成立。
\end{property}

\begin{example}[伯努利随机变量的期望]
若 $X\sim\operatorname{Bernoulli}(p)$，则
\[
    E(X)=0(1-p)+1\cdot p=p.
\]
若 $X_1,\ldots,X_n$ 为相互独立且服从同一伯努利分布的随机变量，则由期望的线性性，
\[
    E\left(\sum_{i=1}^n X_i\right)=\sum_{i=1}^n E(X_i)=np.
\]
\end{example}

\subsection{方差与切比雪夫不等式}
随机变量 $X$ 的方差定义为
\[
    D(X)=\operatorname{Var}(X)=E\bigl[(X-E(X))^2\bigr],
\]
前提是 $E(X^2)<\infty$。方差的平方根
\[
    \sigma_X=\sqrt{D(X)}
\]
称为标准差。方差的等价计算式为
\[
    D(X)=E(X^2)-[E(X)]^2.
\]
离散型与连续型情形分别为
\[
    D(X)=\sum_k(x_k-E(X))^2p_k,\qquad
    D(X)=\int_{-\infty}^{+\infty}(x-E(X))^2f_X(x)\,\mathrm{d}x.
\]

\begin{property}[方差的性质]
对常数 $a,b$，
\[
    D(aX+b)=a^2D(X).
\]
对具有有限二阶矩的 $X,Y$，
\[
    D(X+Y)=D(X)+D(Y)+2\operatorname{Cov}(X,Y).
\]
若 $X,Y$ 独立，则 $\operatorname{Cov}(X,Y)=0$，从而 $D(X+Y)=D(X)+D(Y)$。
\end{property}

\begin{theorem}[切比雪夫不等式]
若 $E(X)=\mu$ 且 $D(X)=\sigma^2<\infty$，则对任意 $\varepsilon>0$，
\[
    P(|X-\mu|\geq\varepsilon)\leq\frac{\sigma^2}{\varepsilon^2}.
\]
\end{theorem}
切比雪夫不等式仅依赖均值和方差，不要求 $X$ 服从某个特定分布。等价地，对任意 $c>0$，
\[
    P(|X-\mu|<c\sigma)\geq 1-\frac{1}{c^2},
\]
其中当右侧为负时，该下界不提供非平凡信息。

\begin{example}[切比雪夫不等式的界]
若 $E(X)=10$、$D(X)=4$，则
\[
    P(|X-10|\geq 4)\leq\frac{4}{4^2}=\frac14,
\]
故 $P(6<X<14)\geq 3/4$。
\end{example}

\subsection{协方差与相关系数}
若 $E(X^2),E(Y^2)<\infty$，则 $X,Y$ 的协方差定义为
\[
    \operatorname{Cov}(X,Y)
    =E\bigl[(X-E(X))(Y-E(Y))\bigr]
    =E(XY)-E(X)E(Y).
\]
协方差的基本性质为
\[
\begin{aligned}
    &\operatorname{Cov}(X,Y)=\operatorname{Cov}(Y,X),\\
    &\operatorname{Cov}(aX+b,cY+d)=ac\operatorname{Cov}(X,Y),\\
    &\operatorname{Cov}(X_1+X_2,Y)
      =\operatorname{Cov}(X_1,Y)+\operatorname{Cov}(X_2,Y).
\end{aligned}
\]
特别地，$D(X)=\operatorname{Cov}(X,X)$。

若 $D(X)>0$ 且 $D(Y)>0$，相关系数定义为
\[
    \rho_{XY}
    =\frac{\operatorname{Cov}(X,Y)}
    {\sqrt{D(X)D(Y)}}.
\]
柯西--施瓦茨不等式给出 $-1\leq\rho_{XY}\leq1$。$\rho_{XY}$ 描述线性相关程度；$\rho_{XY}=0$ 表示不相关。独立且二阶矩存在可推出不相关，但不相关一般不能推出独立。若 $(X,Y)$ 服从二维正态分布，则不相关与独立等价。

\begin{example}[不相关但不独立]
令 $X$ 以相等概率取 $-1,0,1$，并令 $Y=X^2$。则 $E(X)=0$、$E(XY)=E(X^3)=0$，因而
\[
    \operatorname{Cov}(X,Y)=E(XY)-E(X)E(Y)=0.
\]
但 $Y$ 是 $X$ 的确定函数，且
\[
    P(X=0,Y=0)=\frac13
    \ne P(X=0)P(Y=0)=\frac13\cdot\frac13=\frac19,
\]
所以 $X,Y$ 不独立。
\end{example}

\subsection{原点矩、中心矩与协方差矩阵}
若相应期望存在，$X$ 的 $k$ 阶原点矩与 $k$ 阶中心矩分别定义为
\[
    \mu_k'=E(X^k),\qquad
    \mu_k=E\bigl[(X-E(X))^k\bigr].
\]
一阶原点矩为 $E(X)$，二阶中心矩为 $D(X)$。

设随机向量 $\bm{X}=(X_1,\ldots,X_n)^{\mathsf T}$，均值向量为
\[
    \bm{\mu}=E(\bm{X})
    =(E(X_1),\ldots,E(X_n))^{\mathsf T}.
\]
其协方差矩阵定义为
\[
    \bm{\Sigma}
    =E\left[(\bm{X}-\bm{\mu})(\bm{X}-\bm{\mu})^{\mathsf T}\right],
    \qquad
    \Sigma_{ij}=\operatorname{Cov}(X_i,X_j).
\]
协方差矩阵对称且半正定；其对角元素为各分量的方差。

\subsection{多元正态分布}
设 $\bm{X}\in\mathbb{R}^n$，均值向量为 $\bm{\mu}$，协方差矩阵 $\bm{\Sigma}$ 正定。若 $\bm{X}$ 的联合密度为
\[
    f_{\bm{X}}(\bm{x})
    =\frac{1}{(2\pi)^{n/2}|\bm{\Sigma}|^{1/2}}
      \exp\left[
        -\frac12(\bm{x}-\bm{\mu})^{\mathsf T}
        \bm{\Sigma}^{-1}(\bm{x}-\bm{\mu})
      \right],
    \qquad \bm{x}\in\mathbb{R}^n,
\]
则称 $\bm{X}$ 服从 $n$ 维正态分布，记为
\[
    \bm{X}\sim N_n(\bm{\mu},\bm{\Sigma}).
\]
这里 $|\bm{\Sigma}|$ 表示行列式。协方差矩阵正定保证上述密度表达式存在；若矩阵奇异，多元正态分布集中在低维子空间上，不具有关于 $n$ 维 Lebesgue 测度的密度。

\clearpage
\section{大数定律与中心极限定理}
\subsection{依概率收敛}
设 $\{X_n\}$ 为随机变量序列，$X$ 为随机变量。若对任意 $\varepsilon>0$，
\[
    \lim_{n\to\infty}P(|X_n-X|\geq\varepsilon)=0,
\]
则称 $X_n$ 依概率收敛于 $X$，记为 $X_n\xrightarrow{P}X$。若极限对象为常数 $a$，则记为 $X_n\xrightarrow{P}a$。

依概率收敛保持加法与连续变换：若 $X_n\xrightarrow{P}X$、$Y_n\xrightarrow{P}Y$，则
\[
    X_n+Y_n\xrightarrow{P}X+Y,\qquad
    X_nY_n\xrightarrow{P}XY.
\]
若 $g$ 在相应取值处连续，则 $g(X_n)\xrightarrow{P}g(X)$。依概率极限若存在，则以概率为 1 唯一。

\subsection{大数定律}
大数定律描述随机变量平均值在样本量增加时的收敛性质。

\begin{theorem}[切比雪夫大数定律]
设 $X_1,X_2,\ldots$ 相互独立，均有有限方差，且存在常数 $C<\infty$ 使 $D(X_k)\leq C$ 对所有 $k$ 成立。则对任意 $\varepsilon>0$，
\[
    P\left(
        \left|\frac1n\sum_{k=1}^{n}\bigl(X_k-E(X_k)\bigr)\right|
        \geq\varepsilon
    \right)
    \leq \frac{C}{n\varepsilon^2}
    \longrightarrow 0.
\]
因此，
\[
    \frac1n\sum_{k=1}^{n}\bigl(X_k-E(X_k)\bigr)
    \xrightarrow{P}0.
\]
\end{theorem}

\begin{theorem}[独立同分布大数定律]
设 $X_1,X_2,\ldots$ 相互独立同分布，且 $E|X_1|<\infty$，令 $\mu=E(X_1)$。则
\[
    \overline X_n=\frac1n\sum_{k=1}^{n}X_k
    \xrightarrow{P}\mu.
\]
\end{theorem}
若 $X_k$ 为独立同分布的伯努利随机变量，且 $P(X_k=1)=p$，则 $\overline X_n$ 是 $n$ 次试验中成功的频率，从而频率依概率收敛于 $p$。

\subsection{中心极限定理}
大数定律描述标准化前平均值的收敛；中心极限定理描述和或平均值经适当标准化后的分布近似。

\begin{theorem}[独立同分布中心极限定理]
设 $X_1,X_2,\ldots$ 相互独立同分布，且
\[
    E(X_k)=\mu,\qquad D(X_k)=\sigma^2,\qquad 0<\sigma^2<\infty.
\]
则当 $n\to\infty$ 时，
\[
    \frac{\sum_{k=1}^{n}X_k-n\mu}{\sigma\sqrt n}
    =\frac{\sqrt n(\overline X_n-\mu)}{\sigma}
    \xrightarrow{d}N(0,1),
\]
其中 $\xrightarrow{d}$ 表示依分布收敛。
\end{theorem}
对连续极限分布的连续点 $z$，有
\[
    P\left(
        \frac{\sum_{k=1}^{n}X_k-n\mu}{\sigma\sqrt n}\leq z
    \right)
    \longrightarrow \Phi(z),
\]
其中 $\Phi$ 为标准正态分布函数。对有限的 $n$，可用该极限分布构造概率近似，其精度取决于 $n$ 及 $X_k$ 的分布。

\subsection{近似概率计算}
记 $S_n=\sum_{k=1}^{n}X_k$，其中 $X_k$ 独立同分布，均值为 $\mu$、方差为 $\sigma^2>0$。中心极限定理给出
\[
    P(a<S_n\leq b)
    \approx
    \Phi\left(\frac{b-n\mu}{\sigma\sqrt n}\right)
    -\Phi\left(\frac{a-n\mu}{\sigma\sqrt n}\right).
\]
对样本均值 $\overline X_n=S_n/n$，相应地
\[
    P(a<\overline X_n\leq b)
    \approx
    \Phi\left(\frac{(b-\mu)\sqrt n}{\sigma}\right)
    -\Phi\left(\frac{(a-\mu)\sqrt n}{\sigma}\right).
\]
离散型和的概率可使用连续性修正。例如，若 $S_n$ 为整数值，
\[
    P(r\leq S_n\leq s)
    \approx
    \Phi\left(\frac{s+\frac12-n\mu}{\sigma\sqrt n}\right)
    -\Phi\left(\frac{r-\frac12-n\mu}{\sigma\sqrt n}\right).
\]
连续性修正用于将整数区间边界与连续正态近似对应；它不改变精确分布，而是对正态近似作边界调整。

\begin{example}[二项分布的正态近似]
设 $X\sim B(100,0.5)$，则 $E(X)=50$、$D(X)=25$。用中心极限定理并作连续性修正，
\[
\begin{aligned}
    P(45\leq X\leq55)
    &\approx
    \Phi\left(\frac{55.5-50}{5}\right)
    -\Phi\left(\frac{44.5-50}{5}\right)\\
    &=\Phi(1.1)-\Phi(-1.1).
\end{aligned}
\]
\end{example}

\clearpage
\section{样本与抽样分布}
\subsection{总体、样本与简单随机样本}
统计总体是研究对象的全体，样本是从总体中取得的观测值。总体可由分布 $F$ 描述；总体中每个观测对象称为个体，总体所含个体数称为总体容量。

设 $X_1,\ldots,X_n$ 相互独立且服从同一总体分布 $F$，则称其为来自该总体的简单随机样本。观测到的具体数值 $x_1,\ldots,x_n$ 称为样本观测值。用独立同分布样本建模与有限总体中不放回抽样并非同一抽样机制；有限总体不放回抽样的观测一般不独立。

\subsection{直方图与箱线图}
直方图用于展示数值型数据的分布形态。将观测范围划分为互不重叠的组区间 $I_j$，记区间宽度为 $h_j$、落入该区间的观测数为 $n_j$，样本容量为 $n$。若直方图以频率密度为柱高，则第 $j$ 组的柱高为
\[
    d_j=\frac{n_j/n}{h_j}.
\]
此时柱面积 $d_jh_j=n_j/n$ 等于该组的相对频率；各柱面积之和为 $1$。当组距相等时，频数直方图的柱高也可直接用频数表示，但柱面积不再等于相对频率，除非作适当尺度换算。

箱线图以五数概括描述数据的位置与离散程度：最小值、第一四分位数 $Q_1$、中位数 $Q_2$、第三四分位数 $Q_3$ 和最大值。四分位距为
\[
    \operatorname{IQR}=Q_3-Q_1.
\]
按 $1.5\operatorname{IQR}$ 规则，低于 $Q_1-1.5\operatorname{IQR}$ 或高于 $Q_3+1.5\operatorname{IQR}$ 的观测标记为离群点；箱线图的须延伸至区间边界内最远的观测值。该规则用于描述数据，不单独构成关于观测值错误或总体异常的统计结论。

\subsection{统计量与经验分布函数}
统计量是样本的函数，且不含未知总体参数。例如样本均值、样本方差、样本最大值均为统计量：
\[
    \overline X=\frac1n\sum_{i=1}^{n}X_i,\qquad
    S^2=\frac1{n-1}\sum_{i=1}^{n}(X_i-\overline X)^2.
\]
统计量本身是随机变量，其分布称为抽样分布。

样本的经验分布函数定义为
\[
    F_n(x)=\frac1n\sum_{i=1}^{n}\mathbf{1}_{\{X_i\leq x\}},
\]
其中 $\mathbf{1}_{\{X_i\leq x\}}$ 为指示函数。$F_n(x)$ 等于样本观测值中不超过 $x$ 的比例。

\subsection{卡方分布、t 分布与 F 分布}
设 $Z_1,\ldots,Z_\nu$ 相互独立且均服从标准正态分布，则
\[
    U=\sum_{i=1}^{\nu}Z_i^2\sim\chi^2(\nu),
\]
称 $U$ 服从自由度为 $\nu$ 的卡方分布。若 $U_1,U_2$ 独立且分别服从 $\chi^2(\nu_1)$、$\chi^2(\nu_2)$，则
\[
    T=\frac{Z}{\sqrt{U_1/\nu_1}}\sim t(\nu_1),
    \qquad
    F=\frac{U_1/\nu_1}{U_2/\nu_2}\sim F(\nu_1,\nu_2),
\]
其中 $Z\sim N(0,1)$，且 $Z,U_1,U_2$ 相互独立。$t(\nu)$ 表示自由度为 $\nu$ 的 t 分布；$F(\nu_1,\nu_2)$ 的两个自由度依次对应分子与分母中的卡方变量。

卡方分布的自由度可加：若 $U_1,U_2$ 独立且分别服从 $\chi^2(\nu_1)$、$\chi^2(\nu_2)$，则 $U_1+U_2\sim\chi^2(\nu_1+\nu_2)$。t 分布关于 $0$ 对称，且自由度趋于无穷时趋近标准正态分布。若 $F\sim F(\nu_1,\nu_2)$，则 $1/F\sim F(\nu_2,\nu_1)$。

\subsection{正态总体下的重要抽样分布定理}
设 $X_1,\ldots,X_n$ 为来自总体 $N(\mu,\sigma^2)$ 的简单随机样本，定义 $\overline X$ 和 $S^2$ 如上，则
\[
    \overline X\sim N\left(\mu,\frac{\sigma^2}{n}\right),\qquad
    \frac{(n-1)S^2}{\sigma^2}\sim\chi^2(n-1),
\]
且 $\overline X$ 与 $S^2$ 相互独立。因此，
\[
    \frac{\overline X-\mu}{\sigma/\sqrt n}\sim N(0,1)
\]
（$\sigma$ 已知），而
\[
    \frac{\overline X-\mu}{S/\sqrt n}\sim t(n-1)
\]
（$\sigma$ 未知）。

设另有独立样本 $Y_1,\ldots,Y_m$ 来自 $N(\mu_2,\sigma_2^2)$，第一组样本来自 $N(\mu_1,\sigma_1^2)$，两组样本彼此独立。记两组样本均值、样本方差为 $\overline X,S_X^2$ 与 $\overline Y,S_Y^2$。当总体方差已知时，
\[
    \frac{(\overline X-\overline Y)-(\mu_1-\mu_2)}
    {\sqrt{\sigma_1^2/n+\sigma_2^2/m}}
    \sim N(0,1).
\]
当总体方差未知但相等，即 $\sigma_1^2=\sigma_2^2=\sigma^2$，令
\[
    S_p^2=\frac{(n-1)S_X^2+(m-1)S_Y^2}{n+m-2},
\]
则
\[
    \frac{(\overline X-\overline Y)-(\mu_1-\mu_2)}
    {S_p\sqrt{1/n+1/m}}
    \sim t(n+m-2).
\]
若总体方差分别未知且不假定相等，则上述合并方差 t 统计量不适用。

对上述两组正态总体样本，
\[
    \frac{S_X^2/\sigma_1^2}{S_Y^2/\sigma_2^2}
    \sim F(n-1,m-1).
\]
特别地，在 $\sigma_1^2=\sigma_2^2$ 时，$S_X^2/S_Y^2\sim F(n-1,m-1)$。

\clearpage
\section{参数估计}
\subsection{点估计的基本概念}
设总体分布含有未知参数 $\theta$，样本为 $X_1,\ldots,X_n$。由样本构成的统计量 $T=T(X_1,\ldots,X_n)$ 若用于估计 $\theta$，称为 $\theta$ 的估计量。给定样本观测值 $x_1,\ldots,x_n$ 后，$T(x_1,\ldots,x_n)$ 称为估计值。

点估计以一个数值估计参数；区间估计给出包含参数的随机区间及其覆盖概率。估计量在抽样前是随机变量，估计值在样本观测后是确定数值。

\subsection{矩估计法}
设总体含 $k$ 个未知参数 $\theta_1,\ldots,\theta_k$，且所需的前 $k$ 阶原点矩存在。总体原点矩与样本原点矩分别为
\[
    \mu_r'(\bm{\theta})=E(X^r),\qquad
    m_r'=\frac1n\sum_{i=1}^{n}X_i^r,\qquad r=1,\ldots,k.
\]
矩估计法以方程组
\[
    \mu_r'(\bm{\theta})=m_r',\qquad r=1,\ldots,k
\]
确定参数估计量。方程组需有满足参数空间约束的解；解不唯一时，矩估计量不唯一。

\begin{example}[均匀分布参数的矩估计]
设 $X_1,\ldots,X_n$ 来自 $U(a,b)$，其中 $a<b$ 未知。由
\[
    E(X)=\frac{a+b}{2},\qquad
    D(X)=\frac{(b-a)^2}{12},
\]
用样本均值 $\overline X$ 和分母为 $n$ 的样本中心二阶矩
\[
    m_2'-\overline X^2
    =\frac1n\sum_{i=1}^{n}(X_i-\overline X)^2
\]
分别替代总体均值和方差，得矩估计
\[
    \widehat a_{\mathrm{MM}}
    =\overline X-\sqrt{3(m_2'-\overline X^2)},\qquad
    \widehat b_{\mathrm{MM}}
    =\overline X+\sqrt{3(m_2'-\overline X^2)}.
\]
\end{example}

\subsection{极大似然估计法}
设样本的联合分布律或联合密度为 $f(x_1,\ldots,x_n;\theta)$。固定观测值 $x_1,\ldots,x_n$，将其视为参数 $\theta$ 的函数：
\[
    L(\theta)=f(x_1,\ldots,x_n;\theta).
\]
使 $L(\theta)$ 在参数空间内达到最大值的参数值称为极大似然估计值；相应的统计量称为极大似然估计量。若似然函数可微，常改为最大化对数似然 $\ell(\theta)=\log L(\theta)$；解似然方程仍须检验参数空间边界及极值性质。

\begin{example}[伯努利参数的极大似然估计]
设 $X_1,\ldots,X_n$ 独立同分布，且 $P(X_i=1)=p$、$P(X_i=0)=1-p$。令 $s=\sum_{i=1}^{n}x_i$，则似然函数为
\[
    L(p)=p^s(1-p)^{n-s},\qquad 0\leq p\leq1.
\]
当 $0<s<n$ 时，最大化对数似然
\[
    \ell(p)=s\log p+(n-s)\log(1-p)
\]
得到 $\widehat p_{\mathrm{ML}}=s/n$。当 $s=0$ 或 $s=n$ 时，最大值分别在参数空间边界 $p=0$ 或 $p=1$ 处取得，因此同一公式 $\widehat p_{\mathrm{ML}}=\overline X$ 仍成立。
\end{example}

若 $\widehat\theta$ 是 $\theta$ 的极大似然估计，且 $g$ 为一一对应变换，则 $g(\widehat\theta)$ 是 $g(\theta)$ 的极大似然估计。

\subsection{估计量的评选标准}
估计量的性质用于描述其抽样分布表现。
\begin{itemize}
    \item \textbf{无偏性}：若对参数空间中所有 $\theta$，$E_\theta(T)=\theta$，则称 $T$ 为 $\theta$ 的无偏估计量。
    \item \textbf{有效性}：在同一参数的两个无偏估计量中，若 $D_\theta(T_1)\leq D_\theta(T_2)$ 对所有 $\theta$ 成立且至少在某些参数值处严格小于，则 $T_1$ 的方差较小；比较须限定在指定估计量类中。
    \item \textbf{相合性}：若对任意 $\varepsilon>0$，
    \[
        \lim_{n\to\infty}P_\theta(|T_n-\theta|>\varepsilon)=0,
    \]
    则称估计量序列 $T_n$ 依概率相合于 $\theta$。
\end{itemize}
无偏性描述有限样本下的期望，相合性描述样本量趋于无穷时的收敛；二者是不同性质。

\subsection{置信区间与单侧置信区间}
令 $L(\bm X)$、$U(\bm X)$ 为样本统计量，且 $L(\bm X)\leq U(\bm X)$。若对给定参数 $\theta$，
\[
    P_\theta\bigl(L(\bm X)\leq\theta\leq U(\bm X)\bigr)=1-\alpha,
\]
则 $[L(\bm X),U(\bm X)]$ 称为 $\theta$ 的置信水平为 $1-\alpha$ 的置信区间。置信水平是区间构造程序的覆盖概率；参数 $\theta$ 在经典频率学派框架中为固定未知常数。

令 $z_\gamma$ 为标准正态分布的 $\gamma$ 分位数，即 $\Phi(z_\gamma)=\gamma$；令 $t_{\gamma,\nu}$、$\chi^2_{\gamma,\nu}$ 和 $F_{\gamma,\nu_1,\nu_2}$ 分别表示相应分布的 $\gamma$ 分位数。置信区间通常由已知分布的枢轴量反解得到。

单侧置信区间只给出一个方向的界限。例如，若
\[
    P_\theta\bigl(L(\bm X)\leq\theta\bigr)=1-\alpha,
\]
则 $[L(\bm X),+\infty)$ 为置信水平 $1-\alpha$ 的单侧置信区间；上置信限情形类似。

\subsection{正态总体参数的区间估计}
设 $X_1,\ldots,X_n$ 是来自 $N(\mu,\sigma^2)$ 的简单随机样本，样本均值和样本方差分别为 $\overline X$ 与
\[
    S^2=\frac1{n-1}\sum_{i=1}^{n}(X_i-\overline X)^2.
\]
总体均值 $\mu$ 的双侧置信区间为：
\begin{itemize}
    \item $\sigma$ 已知：
    \[
        \left[
        \overline X-z_{1-\alpha/2}\frac{\sigma}{\sqrt n},
        \overline X+z_{1-\alpha/2}\frac{\sigma}{\sqrt n}
        \right];
    \]
    \item $\sigma$ 未知：
    \[
        \left[
        \overline X-t_{1-\alpha/2,n-1}\frac{S}{\sqrt n},
        \overline X+t_{1-\alpha/2,n-1}\frac{S}{\sqrt n}
        \right].
    \]
\end{itemize}
总体方差 $\sigma^2$ 的置信区间（$\mu$ 未知）为
\[
    \left[
        \frac{(n-1)S^2}{\chi^2_{1-\alpha/2,n-1}},
        \frac{(n-1)S^2}{\chi^2_{\alpha/2,n-1}}
    \right].
\]

对两个相互独立的正态总体样本，记样本均值为 $\overline X,\overline Y$，样本方差为 $S_X^2,S_Y^2$，样本容量为 $n,m$。总体均值差 $\mu_1-\mu_2$ 的区间为：
\begin{itemize}
    \item 总体方差 $\sigma_1^2,\sigma_2^2$ 已知：
    \[
        (\overline X-\overline Y)
        \pm z_{1-\alpha/2}
        \sqrt{\frac{\sigma_1^2}{n}+\frac{\sigma_2^2}{m}};
    \]
    \item 总体方差未知且相等，令
    \[
        S_p^2=\frac{(n-1)S_X^2+(m-1)S_Y^2}{n+m-2},
    \]
    则区间为
    \[
        (\overline X-\overline Y)
        \pm t_{1-\alpha/2,n+m-2}S_p\sqrt{\frac1n+\frac1m}.
    \]
\end{itemize}
总体方差比 $\sigma_1^2/\sigma_2^2$ 的置信区间为
\[
    \left[
        \frac{S_X^2/S_Y^2}{F_{1-\alpha/2,n-1,m-1}},
        \frac{S_X^2/S_Y^2}{F_{\alpha/2,n-1,m-1}}
    \right].
\]
以上两样本区间以两组独立正态样本为条件；合并方差均值差区间还要求总体方差相等。

\subsection{0--1 分布参数的区间估计}
设 $X_1,\ldots,X_n$ 独立同分布且 $P(X_i=1)=p$、$P(X_i=0)=1-p$。样本比例
\[
    \widehat p=\overline X=\frac1n\sum_{i=1}^{n}X_i
\]
估计 $p$。由中心极限定理及方差的样本代入，$p$ 的近似双侧置信区间为
\[
    \left[
        \widehat p-z_{1-\alpha/2}
        \sqrt{\frac{\widehat p(1-\widehat p)}{n}},
        \widehat p+z_{1-\alpha/2}
        \sqrt{\frac{\widehat p(1-\widehat p)}{n}}
    \right].
\]
该区间是大样本正态近似，不是有限样本下的精确区间；当样本比例接近 $0$ 或 $1$、样本量较小时，正态近似可能不准确。

\clearpage
\section{假设检验}
\subsection{假设检验的基本思想与步骤}
假设检验利用样本信息判断关于总体分布或参数的假设是否与观测结果相容。原假设记为 $H_0$，备择假设记为 $H_1$。由样本构造检验统计量，并依据其在 $H_0$ 成立时的分布确定拒绝域 $W$。若统计量观测值落入 $W$，则拒绝 $H_0$；否则不拒绝 $H_0$。

检验问题的基本步骤为：
\begin{enumerate}
    \item 根据研究问题提出 $H_0$ 与 $H_1$；
    \item 选定显著性水平 $\alpha$，并指定检验统计量；
    \item 在 $H_0$ 成立的条件下确定统计量分布与拒绝域；
    \item 根据样本计算统计量，并判断其是否落入拒绝域。
\end{enumerate}
不拒绝 $H_0$ 表示样本证据不足以在给定显著性水平下拒绝它，不等同于证明 $H_0$ 为真。

\subsection{显著性水平、检验方向与两类错误}
显著性水平 $\alpha$ 是在 $H_0$ 为真时拒绝 $H_0$ 的概率上界。第一类错误是 $H_0$ 为真而拒绝 $H_0$，其概率记为 $\alpha$（或不超过预先给定的 $\alpha$）；第二类错误是在 $H_1$ 为真时不拒绝 $H_0$，其概率记为 $\beta$，且依赖于具体备择参数值。检验功效为 $1-\beta$。

若备择假设为 $\theta\ne\theta_0$，拒绝域位于统计量分布两侧，称为双侧检验；若 $H_1:\theta>\theta_0$ 或 $H_1:\theta<\theta_0$，则分别为右侧检验或左侧检验。检验方向由备择假设确定。

\subsection{正态总体均值的假设检验}
以下 $z_\gamma$、$t_{\gamma,\nu}$、$\chi^2_{\gamma,\nu}$、$F_{\gamma,\nu_1,\nu_2}$ 分别表示相应分布的 $\gamma$ 分位数。

\paragraph{单个正态总体均值}
设 $X_1,\ldots,X_n$ 来自 $N(\mu,\sigma^2)$，检验 $H_0:\mu=\mu_0$。
\begin{itemize}
    \item $\sigma$ 已知时，
    \[
        Z=\frac{\overline X-\mu_0}{\sigma/\sqrt n}\sim N(0,1)
        \quad(H_0\text{ 成立}).
    \]
    双侧检验 $H_1:\mu\ne\mu_0$ 的拒绝域为 $|Z|>z_{1-\alpha/2}$；右侧检验 $H_1:\mu>\mu_0$ 的拒绝域为 $Z>z_{1-\alpha}$；左侧检验 $H_1:\mu<\mu_0$ 的拒绝域为 $Z<-z_{1-\alpha}$。
    \item $\sigma$ 未知时，令 $S^2=(n-1)^{-1}\sum_{i=1}^{n}(X_i-\overline X)^2$，则
    \[
        T=\frac{\overline X-\mu_0}{S/\sqrt n}\sim t(n-1)
        \quad(H_0\text{ 成立}).
    \]
    双侧检验拒绝域为 $|T|>t_{1-\alpha/2,n-1}$；右侧检验拒绝域为 $T>t_{1-\alpha,n-1}$；左侧检验拒绝域为 $T<-t_{1-\alpha,n-1}$。
\end{itemize}

\begin{example}[已知方差时的均值双侧检验]
设总体为 $N(\mu,4)$，抽取容量 $n=16$ 的简单随机样本，检验 $H_0:\mu=10$ 对 $H_1:\mu\ne10$，显著性水平为 $\alpha=0.05$。此时
\[
    Z=\frac{\overline X-10}{2/4}.
\]
由于 $z_{0.975}\approx1.96$，拒绝域为 $|Z|>1.96$，等价于 $\overline X<9.02$ 或 $\overline X>10.98$。
\end{example}

设另有容量为 $m$ 的独立正态样本 $Y_1,\ldots,Y_m$，两组总体均值为 $\mu_1,\mu_2$，检验 $H_0:\mu_1-\mu_2=\Delta_0$。
\begin{itemize}
    \item 若总体方差已知，
    \[
        Z=\frac{(\overline X-\overline Y)-\Delta_0}
        {\sqrt{\sigma_1^2/n+\sigma_2^2/m}}
        \sim N(0,1)
        \quad(H_0\text{ 成立}).
    \]
    \item 若总体方差未知且相等，定义
    \[
        S_p^2=\frac{(n-1)S_X^2+(m-1)S_Y^2}{n+m-2},
    \]
    则
    \[
        T=\frac{(\overline X-\overline Y)-\Delta_0}
        {S_p\sqrt{1/n+1/m}}
        \sim t(n+m-2)
        \quad(H_0\text{ 成立}).
    \]
\end{itemize}
两侧检验、右侧检验和左侧检验分别依据相应标准正态或 t 分布的双尾、右尾和左尾分位数确定拒绝域。

\subsection{正态总体方差的假设检验}
设 $X_1,\ldots,X_n$ 为来自 $N(\mu,\sigma^2)$ 的简单随机样本，检验 $H_0:\sigma^2=\sigma_0^2$，其中 $\sigma_0^2>0$。令
\[
    Q=\frac{(n-1)S^2}{\sigma_0^2}.
\]
在 $H_0$ 成立时，$Q\sim\chi^2(n-1)$。双侧备择 $H_1:\sigma^2\ne\sigma_0^2$ 的拒绝域为
\[
    Q<\chi^2_{\alpha/2,n-1}
    \quad\text{或}\quad
    Q>\chi^2_{1-\alpha/2,n-1}.
\]
右侧备择 $H_1:\sigma^2>\sigma_0^2$ 的拒绝域为 $Q>\chi^2_{1-\alpha,n-1}$；左侧备择 $H_1:\sigma^2<\sigma_0^2$ 的拒绝域为 $Q<\chi^2_{\alpha,n-1}$。

设两组独立正态总体样本容量分别为 $n,m$，样本方差为 $S_X^2,S_Y^2$。检验
\[
    H_0:\frac{\sigma_1^2}{\sigma_2^2}=r_0,\qquad r_0>0.
\]
统计量
\[
    F=\frac{S_X^2/S_Y^2}{r_0}
\]
在 $H_0$ 成立时服从 $F(n-1,m-1)$。对于等方差假设 $H_0:\sigma_1^2=\sigma_2^2$，取 $r_0=1$，则 $F=S_X^2/S_Y^2$。双侧检验拒绝域为
\[
    F<F_{\alpha/2,n-1,m-1}
    \quad\text{或}\quad
    F>F_{1-\alpha/2,n-1,m-1}.
\]
右侧备择 $\sigma_1^2/\sigma_2^2>r_0$ 的拒绝域为 $F>F_{1-\alpha,n-1,m-1}$；左侧备择 $\sigma_1^2/\sigma_2^2<r_0$ 的拒绝域为 $F<F_{\alpha,n-1,m-1}$。

\end{document}
